\color{unidarkgreen}{\subsection{Логика отключения присоединений}}
\color{black}

\begin{enumerate}[label=\arabic{section}.\arabic{subsection}.\arabic{enumi}, labelsep=4pt, leftmargin=0pt, itemindent=57pt, itemsep=0pt, parsep=5pt] % если что можно крутить itemsep

\item
Для формирования сигналов отключения и запрета АПВ присоединений предусмотрена возможность оперативного вывода действия защит на присоединение посредством внешнего сигнала \textit{«Вывод действия на Вn»}.

\item
Сигнал отключения присоединения (\textit{«Отключение Вn»}) формируется при отключении СШ, а также при отключении в результате опробования присоединения.

\item
Сигнал аварийного отключения присоединения (\textit{«Откл. аварийно Вn»}) дополнительно формируется при действии УРОВ на себя.

\item
Для каждого присоединения предусмотрена возможность формирования как общего, так и избирательного запрета АПВ.

\item
Общий запрет АПВ присоединения формируется по факту наличия сигнала запрет АПВ СШ. 

\item
Избирательный запрет АПВ присоединения формируется по факту срабатывания ДЗО.

\medskip

\begin{figure}[H]
\centering
\includegraphics[width=0.8\textwidth,height=0.8\textheight,keepaspectratio]{\fbpath/761.9901 Логика отключения (ДЗО)/img/Log_otkl_Q1.pdf}
\captionof{figure}{Функциональный блок логики отключения первого присоединения}\label{pict:Log_otkl_Q1}
\end{figure}

\medskip

\item
Уставки логики отключения присоединений представлены в таблице \ref{desc:tbl1}.

\small
\begin{longtable}{|>{\centering\arraybackslash}m{5.3cm}|>{\centering\arraybackslash}m{3.3cm}|>{\centering\arraybackslash}m{4.2cm}|>{\centering\arraybackslash}m{1.8cm}|>{\centering\arraybackslash}m{1cm}|}
% Первая страница: полный заголовок
\caption{Перечень уставок логики отключения присоединений\hfill\vspace{-0.5\baselineskip}}\label{desc:tbl1}\\
\hline
\rowcolor{gray!30}
Параметр & Условное обозначение на схеме & Значение/ Диапазон & Единица измерения & Шаг \\
\hline
\endfirsthead

% Промежуточные страницы: «Продолжение таблицы»
\caption*{\hspace{3pt}\emph{Продолжение таблицы \ref{desc:tbl1}}\hfill\ }\vspace{-0.5\baselineskip} \\
\hline
\rowcolor{gray!30}
Параметр & Условное обозначение на схеме & Значение/ Диапазон & Единица измерения & Шаг \\
\hline
\endhead

%===m>ДЗО|-
%===m
\end{longtable}
\normalsize	


\end{enumerate}